
\subsection{Sylvesterの判定法2}

このsubsectionでは$M \in \M_{n+1}(\R)$とする．

\begin{definition}
  \label{def:lastCol}
  \lean{lastCol}
  \leanok
  ベクトル$b := (b_1, \ldots, b_n) \in \R^n$を$b_i = M_{i\ n+1}$で定める．
  すなわち，$b$は$M$の最後の列から一番下の行を除いた縦ベクトルである．
\end{definition}

\begin{definition}
  \label{def:lastEntry}
  \lean{lastEntry}
  \leanok
  定数$c \in \R$を$c = M_{n+1\ n+1}$で定める．
  すなわち，$c$は$M$の一番右下の成分である．
\end{definition}

これらの定義を踏まえて，以下
\begin{equation}
  M = \left(
    \begin{array}{ccc|c}
      &   & & \\
      & A & & b \\
      &   & & \\ \hline
      & \transpose{b} & & c
    \end{array}
  \right)
\end{equation}
とする．

\begin{definition}
  \label{def:schurComplement}
  \lean{schurComplement}
  \leanok
  次の定数$s \in \R$を$M$の\textbf{Schur complement}という：
  \begin{equation}
    s = c - \transpose{b} A^{-1} b.
  \end{equation}
\end{definition}

\begin{lemma}
  \label{lem:quadratic_form_decomp}
  \lean{quadratic_form_decomp}
  \leanok
  $M$はエルミートであるとする．
  任意のベクトル$v := (x_1, \ldots, x_n, t) \in \R^{n+1}$に対し，
  \begin{equation}
    \transpose{v} M v
    = \transpose{x} A x + 2t(\transpose{b}x) + ct^2
  \end{equation}
  が成り立つ．
\end{lemma}

\begin{lemma}
  \label{lem:quadratic_form_complete_square}
  \lean{quadratic_form_complete_square}
  \leanok
  \uses{lem:quadratic_form_decomp}
  $M$はエルミートであるとする．
  上の補題~\ref{lem:quadratic_form_decomp}を平方完成すると以下のようになる：
  \begin{equation}
    \transpose{v} M v
    = \transpose{(x + t A^{-1} b)} A (x + t A^{-1} b) + st^2
  \end{equation}
\end{lemma}

\begin{lemma}
  \label{lem:schur_complement_pos}
  \lean{schur_complement_pos}
  \leanok
  $M$はエルミートであるとする．
  また，$A$は正定値，$\det M > 0$であるとする．
  このとき，$M$のSchur complement $s$は正である：
  \begin{equation}
    0 < s = c - \transpose{b} A^{-1} b.
  \end{equation}
\end{lemma}

\begin{proof}
  $\det M = (\det A) \cdot s$である．
  また，仮定より$A$は正定値，すなわち$\det A > 0$であり，$\det M > 0$でもあるので，$s > 0$が従う．
\end{proof}

\begin{lemma}
  \label{lem:posDef_dotProduct_nonneg}
  \lean{posDef_dotProduct_nonneg}
  \leanok
  $A \in \M_k(\R)$を正定値行列とする．
  また，$x \in \R^k$とする．
  このとき，$\transpose{x} A x \ge 0$である．
\end{lemma}

\begin{theorem}
  \label{thm:posDef_of_leading_submatrix_posDef_and_det_pos}
  \lean{posDef_of_leading_submatrix_posDef_and_det_pos}
  \leanok
  \uses{lem:schur_complement_pos, lem:posDef_dotProduct_nonneg}
  $M$はエルミートであるとする．
  また，$A$は正定値，$\det M > 0$であるとする．
  このとき，$M$は正定値である．
\end{theorem}

\begin{proof}
  任意の$0 \ne v := (x_1, \ldots, x_n, t) \in \R^{n+1}$に対し，$\transpose{v} M v > 0$を示す．

  $t = 0$のとき，$A$の正定値性から
  \begin{equation}
    \transpose{v} M v = \transpose{x} A x > 0
  \end{equation}
  であるからよい．

  $t \ne 0$のとき，補題~\ref{lem:quadratic_form_complete_square}より
  \begin{equation}
    \transpose{v} M v = \transpose{(x + t A^{-1} b)} A (x + t A^{-1} b) + st^2
  \end{equation}
  である．
  右辺の第1項目は$A$の正定値性より非負，第2項目は補題~\ref{lem:schur_complement_pos}より$s > 0$であるから正である．
  よって，この場合でも$\transpose{v} M v > 0$である．
\end{proof}

\begin{lemma}
  \label{lem:leadingSubmatrix_leadingSubmatrix}
  \lean{leadingSubmatrix_leadingSubmatrix}
  \leanok
  $j \le k \le n$とする．
  $M$の$k$次主小行列の$j$次主小行列は，$M$の$j$次主小行列である：
  \begin{equation}
    (M_k)_j = M_j.
  \end{equation}
\end{lemma}

\begin{theorem}
  \label{thm:leading_minors_pos_implies_posDef}
  \lean{leading_minors_pos_implies_posDef}
  \leanok
  \uses{thm:posDef_of_leading_submatrix_posDef_and_det_pos, lem:IsHermitian.submatrix_injective, lem:leadingSubmatrix_leadingSubmatrix}
  $M \in \M_n(\R)$はエルミートであるとする．
  また，任意の$k \le n$に対し，$0 < k$ならば$0 < M_k$であるとする．
  このとき，$M$は正定値である．
\end{theorem}